\subsection{Heterogeneous Ensemble}

To test the effect of architectural diversity, we construct a heterogeneous ensemble from four backbones with distinct biases, all fine-tuned on the same ImageNet-subset training data.

These backbones are summarized in Table~\ref{tab:hetero-ensemble}.
They were specifically chosen to maximize inductive bias diversity: the CNNs provide complementary local feature extraction (texture, details, scaling), while the Vision Transformer (ViT) adds global reasoning.
Identical training ensures differences arise from architecture alone.

\begin{table}[!htbp]
  \caption{Backbone architectures comprising the heterogeneous ensemble. Each model is initialized with ImageNet-pretrained weights and fine-tuned on the same training data to isolate architectural inductive bias effects.}
  \label{tab:hetero-ensemble}
  \centering
  \begin{tabularx}{\linewidth}{
    l
    >{\hsize=1.2\hsize}L
    >{\hsize=0.8\hsize}L
    }
    
    \toprule
    Backbone        & Key characteristics                           & Reference \\
    \midrule
    ResNet-34       & strong local texture bias, stable baseline    & \cite{He_2016_CVPR} \\
    DenseNet-121    & Dense connectivity for feature reuse          & \cite{Huang_2017_CVPR} \\
    EfficientNet-B0 & Parameter-efficient, balanced scaling         & \cite{pmlr-v97-tan19a} \\
    ViT-B/16        & Global self-attention, shape-focused          & \cite{dosovitskiyImageWorth16x162020} \\
    \bottomrule
  \end{tabularx}
\end{table}

Predictions are aggregated via simple averaging:
\[
\hat{y} = \frac{1}{K} \sum_{k=1}^{K} p_k(\mathbf{x})
\]
where $p_k(\mathbf{x})$ is the softmax output of expert $k$, and $K=4$. This uniform weighting provides a baseline before introducing learned routing in Section ~\ref{sec:moe}.

Prior work suggests that architectural diversity can induce complementary failure modes under distribution shift~\citep{wuExploringModelLearning2023}. We evaluate this hypothesis empirically.